\documentclass{iopconfser}

% Common packages
\usepackage{amsmath}
\usepackage{graphicx}
\usepackage{booktabs}
\usepackage{cite}

\begin{document}

\title{Data-Intensive Methods in [Experiment or Observatory]:
A Review of [Physics Measurement]}

\author{Author Name$^{1}$, Author Name$^{2}$ and Author Name$^{1,2}$}

\affil{$^1$Department, Institution, City, Country}
\affil{$^2$Department, Institution, City, Country}

\email{author@institution.org}

\begin{abstract}
Briefly describe the physics question, methods, main result,
and conclusion of your review.
\end{abstract}

\section{Motivation}

\subsection{Physics question and significance}
What does the experiment measure? Explain the hypothesis
being tested and why the measurement matters~\cite{paper}.

\section{Method}

\subsection{Experimental setup}
Describe the apparatus and its main components.
Explain how it detects the physical signal.

\begin{figure}[htbp]
    \centering
    % Replace the box with the image command below when ready.
    \fbox{\parbox[c][4cm][c]{0.8\linewidth}{
        \centering Labelled detector or apparatus diagram
    }}
    % \includegraphics[width=0.8\linewidth]{detector.png}
    \caption{Describe the apparatus. Add the image source
    and attribution here~\cite{paper}.}
    \label{fig:detector}
\end{figure}

\subsection{Data acquisition and processing}
Describe the journey from raw data to the final dataset.
Include numerical values and references for:
\begin{itemize}
    \item Raw event rate or data streaming rate.
    \item Trigger or filtering stages and their input/output rates.
    \item Stored data rate and annual or total recorded volume.
    \item Offline processing, computing system, and total storage.
    \item Simulation steps and their cost relative to data processing.
\end{itemize}

\subsection{Data analysis}
Explain data selection, key observables, background or noise
estimation, and signal extraction. Describe any machine learning
or other computational techniques used.

An example equation is
\begin{equation}
    y = mx + c,
    \label{eq:example}
\end{equation}
where \(m\) is the slope and \(c\) is the intercept.
Replace equation~\eqref{eq:example} with a relevant equation.

\subsection{Systematic uncertainties}
Identify the main uncertainties and explain how they are
evaluated and included in the analysis.

\section{Results}

\subsection{Statistical method}
Describe the statistical framework and test statistic.
Explain how uncertainties, significance, or limits are obtained.

\subsection{Main result}
Report the main result with its statistical and systematic
uncertainties, where available. Explain the key features of
figure~\ref{fig:result}.

\begin{figure}[htbp]
    \centering
    % Replace the box with the image command below when ready.
    \fbox{\parbox[c][4cm][c]{0.8\linewidth}{
        \centering Key result figure from the publication
    }}
    % \includegraphics[width=0.8\linewidth]{result.png}
    \caption{Describe the result. Add the publication,
    original figure number, and attribution here~\cite{paper}.}
    \label{fig:result}
\end{figure}

\subsection{Discussion and critique}
Explain what the result means. Discuss the strengths,
limitations, and role of data-intensive methods.

\section{Conclusion}
Summarise the physics finding and how data-intensive methods
made the measurement possible. Suggest one future improvement.

\section*{Author Contributions}
All authors contributed equally to the literature review,
writing, and revision of this manuscript.
All authors read and approved the final manuscript.

\section*{Acknowledgments}
Acknowledge any assistance or support, or remove this section.

\bibliographystyle{unsrt}
\bibliography{references}

\end{document}